\appendix
\part{Anhang}

\newcommand{\lit}[3]{\vspace{5pt}\noindent\textbf{[#1]} #2\par\nopagebreak\vspace{1pt}{\small\itshape\color{grau} #3}\par\vspace{2pt}}

\chapter{Formelsammlung und Übersicht}

\section{Wachstum}

\begin{merksatz}
\boxtitel{Eine Population}
\begin{center}
{\small
\renewcommand{\arraystretch}{1.62}
\begin{tabular}{@{}P{4.4cm}P{4.2cm}@{}}
\toprule
$\dfrac{\dd N}{\dd t} = rN$ & exponentiell (1.1)\\
$N(t) = N_0e^{rt}$ & Satz 1.2\\
$T_2 = \dfrac{\ln 2}{r}$ & Satz 1.3\\
$T_{1/2} = \dfrac{\ln 2}{k}$ & Zerfall\\
$\dfrac{\dd N}{\dd t} = rN\left(1-\dfrac{N}{K}\right)$ & logistisch (2.1)\\
$N(t) = \dfrac{K}{1+Be^{-rt}}$, $B = \dfrac{K-N_0}{N_0}$ & Satz 2.2\\
Wendepunkt bei $K/2$; $\dot N_{\max} = rK/4$ & Satz 2.3\\
$t_p = \dfrac{1}{r}\ln\dfrac{pB}{1-p}$ & Zeit bis Anteil $p$\\
$\dfrac{\dd N}{\dd t} = rN\ln\dfrac{K}{N}$ & Gompertz (13.1)\\
$\ln\dfrac{K}{N} = \ln\dfrac{K}{N_0}e^{-rt}$ & Satz 13.2\\
Wendepunkt bei $K/e$; $\dot N_{\max} = rK/e$ & Satz 13.3\\
$x_{n+1} = ax_n(1-x_n)$ & logistische Abbildung (3.1)\\
$x^{*} = 0$ und $1-1/a$ & Satz 3.3\\
stabil $\iff |f'(x^{*})|<1$ & Satz 3.4\\
$\delta = 4{,}6692\ldots$ & Feigenbaum\\
\bottomrule
\end{tabular}}
\end{center}
\end{merksatz}

\section{Wechselwirkung}

\begin{merksatz}
\boxtitel{Zwei Populationen}
\begin{center}
{\small
\renewcommand{\arraystretch}{1.62}
\begin{tabular}{@{}P{4.6cm}P{4.0cm}@{}}
\toprule
$\dot x = x(a-by)$, $\dot y = y(-c+dx)$ & Lotka-Volterra (4.1)\\
$(x^{*},y^{*}) = (c/d,\ a/b)$ & Satz 4.2, überkreuz\\
$V = dx - c\ln x + by - a\ln y$ & Erhaltungsgröße (4.3)\\
$\overline x = c/d$, $\overline y = a/b$ & Satz 4.4\\
$T = 2\pi/\sqrt{ac}$ & Periode\\
$\overline x = \dfrac{c+h}{d}$, $\overline y = \dfrac{a-h}{b}$ & Volterras Prinzip\\
$\dot x = r_1x\left(1-\dfrac{x+\alpha_{12}y}{K_1}\right)$ & Konkurrenz (5.1)\\
$\alpha_{12} < \dfrac{K_1}{K_2}$ und $\alpha_{21} < \dfrac{K_2}{K_1}$ & Koexistenz (5.3)\\
$x^{*} = \dfrac{K_1-\alpha_{12}K_2}{1-\alpha_{12}\alpha_{21}}$ & Satz 5.4\\
\bottomrule
\end{tabular}}
\end{center}
\end{merksatz}

\begin{merksatz}
\boxtitel{Stabilität}
\begin{center}
{\small
\renewcommand{\arraystretch}{1.62}
\begin{tabular}{@{}P{4.4cm}P{4.2cm}@{}}
\toprule
$J = \begin{pmatrix} f_x & f_y\\ g_x & g_y\end{pmatrix}$ & Definition 6.1\\
$T = \operatorname{spur}J$, $D = \det J$ & \\
$\lambda_{1,2} = \dfrac{T\pm\sqrt{T^2-4D}}{2}$ & \\
stabil $\iff T<0$ und $D>0$ & Satz 6.3\\
$D<0$ & Sattel\\
$T^2-4D<0$ & Strudel; $>0$ Knoten\\
$T=0$, $D>0$ & Zentrum (unentschieden)\\
$T=0$ überquert & Hopf-Bifurkation\\
\bottomrule
\end{tabular}}
\end{center}
\end{merksatz}

\section{Ausbreitung}

\begin{merksatz}
\boxtitel{Epidemien}
\begin{center}
{\small
\renewcommand{\arraystretch}{1.62}
\begin{tabular}{@{}P{4.4cm}P{4.2cm}@{}}
\toprule
$\dot S = -\beta SI/N$, $\dot I = \beta SI/N - \gamma I$ & SIR (7.1)\\
$\Rnull = \beta/\gamma$ & Satz 7.2\\
$\mathcal{R}_{\text{eff}} = \Rnull\,S/N$ & \\
$\Rnull = 1 + r/\gamma$ & aus der Wachstumsrate\\
$\dfrac{I_{\max}}{N} = 1-\dfrac{1+\ln\Rnull}{\Rnull}$ & Satz 7.3\\
$S$ am Gipfel $= N/\Rnull$ & \\
$s_\infty = e^{-\Rnull(1-s_\infty)}$ & Satz 7.4\\
$p_c = 1-1/\Rnull$ & Satz 8.1\\
$p > p_c/\varepsilon$ & mit Wirksamkeit\\
$S^{*}/N = 1/\Rnull$ & endemisch\\
$A \approx L/(\Rnull-1)$ & Satz 8.2\\
$\Rnull = T\left(\dfrac{\langle k^2\rangle}{\langle k\rangle}-1\right)$ & Satz 9.2, Netz\\
$\dfrac{\langle k^2\rangle}{\langle k\rangle} \ge \langle k\rangle$ & Freundschaftsparadoxon\\
\bottomrule
\end{tabular}}
\end{center}
\end{merksatz}

\section{Der Körper}

\begin{merksatz}
\boxtitel{Pharmakokinetik und Kinetik}
\begin{center}
{\small
\renewcommand{\arraystretch}{1.62}
\begin{tabular}{@{}P{4.4cm}P{4.2cm}@{}}
\toprule
$\mathrm{CL} = k\,V$, $t_{1/2} = \dfrac{\ln2\,V}{\mathrm{CL}}$ & Definition 10.1\\
$c(t) = c_\infty(1-e^{-kt})$, $c_\infty = R/\mathrm{CL}$ & Satz 10.2\\
$1-(1/2)^n$ nach $n$ Halbwertszeiten & Satz 10.3\\
$c_{\max} = \dfrac{D/V}{1-r}$, $r = e^{-k\tau}$ & Satz 10.4\\
$c_{\max}/c_{\min} = 2^{\tau/t_{1/2}}$ & \\
$D_{\text{Auf}} = c_{\text{Ziel}}\,V$ & aus $V$, nicht $\mathrm{CL}$\\
$\mathrm{AUC} = F\,D/\mathrm{CL}$ & \\
$t_{\max} = \dfrac{\ln(k_a/k)}{k_a-k}$ & Bateman\\
$v = \dfrac{V_{\max}[S]}{K_M+[S]}$ & Satz 11.1\\
$K_M = \dfrac{k_{-1}+k_2}{k_1}$ & keine Bindungsstärke\\
$K_{M,\text{eff}} = K_M(1+[I]/K_i)$ & kompetitiv\\
$\theta = \dfrac{[L]^n}{K^n+[L]^n}$ & Hill\\
$[L]_{90}/[L]_{10} = 81^{1/n}$ & Steilheit\\
$E = \dfrac{61{,}5}{z}\lg\dfrac{c_a}{c_i}$ mV & Nernst\\
$\tau = R_mC_m$ & Satz 12.1\\
$v^{*2} = 1-\varepsilon b$ & Satz 12.3, Hopf\\
\bottomrule
\end{tabular}}
\end{center}
\end{merksatz}

\section{Krankheit}

\begin{merksatz}
\boxtitel{Tumor und Resistenz}
\begin{center}
{\small
\renewcommand{\arraystretch}{1.62}
\begin{tabular}{@{}P{4.4cm}P{4.2cm}@{}}
\toprule
Log-kill: fester Anteil je Gabe & 13.4\\
$\Pr[\text{Heilung}] = e^{-\mu}$ & Poisson, 13.5\\
$\Pr[\text{keine Resistenz}] = e^{-\mu N}$ & Satz 14.1\\
$N_c = 1/\mu$ & kritische Zellzahl\\
$\mu_{12} = \mu_1\mu_2$ & nur bei Unabhängigkeit\\
$\operatorname{Var}/\text{Mittel} \gg 1$ & Satz 14.2, vorbestehend\\
$\Pr[\text{keine Mutante}] = e^{-m}$ & Nullklassen-Methode\\
Faktor $1-L/K$ & Verlangsamung bei Eindämmung\\
\bottomrule
\end{tabular}}
\end{center}
\end{merksatz}

\section{Die Landkarte des Bandes}

\begin{merksatz}
\boxtitel{Welcher Satz wo --- und worauf er beruht}
\begin{center}
{\small
\renewcommand{\arraystretch}{1.42}
\begin{tabular}{@{}P{3.6cm}P{1.2cm}P{3.4cm}@{}}
\toprule
\textbf{Ergebnis} & \textbf{Kap.} & \textbf{Beruht auf}\\
\midrule
$N = N_0e^{rt}$ & $1$ & Trennung der Veränderlichen\\
$T_2$ bestandsunabhängig & $1$ & $e^{rt}$ kürzt sich\\
Logistische Lösung & $2$ & Partialbruch\-zerlegung\\
Wendepunkt $K/2$ & $2$ & zweite Ableitung\\
Symmetrie & $2$ & $B$ fällt heraus\\
Fixpunktstabilität & $3$ & Tangentennäherung\\
Perioden\-verdopplung & $3$ & $|f'|$ überschreitet $1$\\
Gleichgewicht überkreuz & $4$ & beide Gleichungen null\\
Erhaltungsgröße & $4$ & Terme heben sich auf\\
\textbf{Volterras Prinzip} & $4$ & $\ln x$ über eine Periode\\
Koexistenz\-bedingung & $5$ & Lage der Nullklinen\\
Satz 5.3 bewiesen & $6$ & $D > 0$\\
Lotka-Volterra ist Zentrum & $6$ & $T=0$\\
$\Rnull = \beta/\gamma$ & $7$ & $\dd I/\dd t > 0$\\
Gipfelhöhe & $7$ & $\dd I/\dd S$ integrieren\\
Endgröße & $7$ & dieselbe Gleichung bei $I=0$\\
Zwei Lösungen & $7$ & Konvexität\\
Herdenschwelle & $8$ & $S/N = 1-p$\\
Altersverschiebung & $8$ & $A = L/(\Rnull-1)$\\
Netz-$\Rnull$ & $9$ & Auswahl über Kanten\\
Anflutung & $10$ & wie Satz 10.2\\
Michaelis-Menten & $11$ & Quasi-Stationarität\\
Hill-Steilheit & $11$ & $\theta/(1-\theta)$\\
Erregbarkeit & $12$ & kubische Nullkline\\
Gompertz-Lösung & $13$ & Substitution $\ln(K/N)$\\
Goldie-Coldman & $14$ & Poisson\\
Luria-Delbrück & $14$ & Varianz gegen Mittelwert\\
\bottomrule
\end{tabular}}
\end{center}
\end{merksatz}

\begin{merksatz}
\boxtitel{Was in diesem Band \kursiv{nicht} bewiesen wurde}
\begin{center}
{\small
\renewcommand{\arraystretch}{1.5}
\begin{tabular}{@{}P{3.4cm}P{1.4cm}P{3.4cm}@{}}
\toprule
\textbf{Ergebnis} & \textbf{Kap.} & \textbf{Wo es steht}\\
\midrule
Trennung der Veränderlichen & $1$ & Band 14, Teil~I\\
Taylor in zwei Veränderlichen & $6$ & Band 3, Teil~V\\
Eigenwerte und Lösungen & $6$ & Band 2, Teil~IV\\
Bateman-Funktion & $10$ & Band 14, Teil~II\\
Feigenbaum-Universalität & $3$ & Strogatz\zit{27}\\
Luria-Delbrück-Verteilung & $14$ & Otto \& Day\zit{24}\\
\bottomrule
\end{tabular}}
\end{center}

\begin{center}
\begin{minipage}{0.88\textwidth}
\kursiv{Sechs Auslassungen, jede im Text benannt, mit Begründung und Fundstelle. Das ist die einzige zulässige Form, die Regel dieses Werkes zu brechen: offen und mit Quellenangabe.}
\end{minipage}
\end{center}
\end{merksatz}

\section{Zeittafel}

\begin{merksatz}
\boxtitel{Von Malthus bis zur adaptiven Therapie}
\begin{center}
{\small
\renewcommand{\arraystretch}{1.36}
\begin{tabular}{@{}P{1.4cm}P{2.8cm}P{4.0cm}@{}}
\toprule
\textbf{Jahr} & \textbf{Person} & \textbf{Ereignis}\\
\midrule
$1798$ & Malthus & geometrisches Wachstum\\
$1825$ & Gompertz & Sterbegesetz\\
$1838$ & Verhulst & logistische Gleichung\\
$1859$ & Darwin & Auslese aus Überproduktion\\
$1890$ & Poincaré & Phasenporträt, Chaos\\
$1892$ & Ljapunow & Stabilitätssatz\\
$1902$ & Nernst & Gleichgewichtsspannung\\
$1910$ & Hill & Sättigungsgleichung\\
$1911$ & Ross & Schwellenwert bei Malaria\\
$1913$ & Michaelis, Menten & Enzymkinetik\\
$1920$ & Lotka & chemische Schwingungen\\
$1925$ & Briggs, Haldane & Quasi-Stationarität\\
$1926$ & Volterra & Räuber-Beute, Prinzip\\
$1927$ & Kermack, McKendrick & SIR-Modell\\
$1932$ & Widmark & Alkoholabbau\\
$1934$ & Gause & Konkurrenzausschluss\\
$1943$ & Luria, Delbrück & Schwankungsversuch\\
$1952$ & Hodgkin, Huxley & Aktionspotential\\
$1953$ & Dost & Begriff Pharmakokinetik\\
$1961$ & FitzHugh & Reduktion auf zwei Gleichungen\\
$1963$ & Lorenz & Wettermodell, Chaos\\
$1964$ & Laird & Gompertz für Tumoren\\
$1971$ & Rosenzweig & Anreicherungsparadoxon\\
$1972$ & May & Vielfalt und Stabilität\\
$1975$ & Feigenbaum & die Konstante\\
$1976$ & May & einfache Modelle, komplexe Dynamik\\
$1979$ & Goldie, Coldman & Resistenz und Tumorgröße\\
$1991$ & Feld & Freundschaftsparadoxon\\
$1999$ & Barabási, Albert & Netze mit Potenzgesetz\\
$2001$ & Pastor-Satorras & Ausbreitung ohne Schwelle\\
$2005$ & Lloyd-Smith & Überstreuung\\
$2009$ & Gatenby & adaptive Therapie\\
\bottomrule
\end{tabular}}
\end{center}
\end{merksatz}

\section{Glossar}

\begin{merksatz}
\boxtitel{Begriffe von A bis Z}
\begin{center}
{\small
\renewcommand{\arraystretch}{1.3}
\begin{longtable}{@{}P{3.6cm}P{5.0cm}@{}}
\toprule
\textbf{Begriff} & \textbf{Bedeutung}\\
\midrule
\endhead
Allee-Effekt & negative Rate bei kleinem Bestand (2)\\
Anreicherungs\-paradoxon & reichere Umwelt, instabiler (6)\\
AUC & Fläche unter der Spiegelkurve (10)\\
Bateman-Funktion & Spiegel nach oraler Gabe (10)\\
Bifurkation & Verhaltenswechsel bei Parameteränderung (6)\\
Bioverfügbarkeit & Anteil, der ins Blut gelangt (10)\\
Chronaxie & $\tau\ln2$, Erregbarkeitsmaß (12)\\
Clearance & geklärtes Volumen je Zeit (10)\\
Depolarisations\-block & Unerregbarkeit bei starkem Reiz (12)\\
Fitnesskosten & Wachstumsnachteil der Resistenz (1, 14)\\
Fixpunkt & $f(x^{*}) = x^{*}$ (3.2)\\
Freundschafts\-paradoxon & $\langle k^2\rangle/\langle k\rangle \ge \langle k\rangle$ (9)\\
Gompertz-Modell & logarithmische Bremsung (13.1)\\
Gradverteilung & Anteil mit je $k$ Kontakten (9.1)\\
Grenzzyklus & isolierte geschlossene Bahn (4)\\
Herdenschwelle & $1-1/\Rnull$ (8.1)\\
Hill-Koeffizient & Maß der Kooperativität (11)\\
Hopf-Bifurkation & aus Ruhe wird Schwingung (6, 12)\\
Identifizierbarkeit & Parameter aus Daten bestimmbar (15)\\
Jacobi-Matrix & Tangente in zwei Veränderlichen (6.1)\\
Kapazitätsgrenze & $K$ im logistischen Modell (2.1)\\
Kollateral\-sensitivität & Resistenz erzeugt andere Schwäche (14)\\
Kolonisations\-resistenz & Vorrangeffekt als Schutz (5)\\
Konkurrenz\-ausschluss & gleiche Nische, keine Koexistenz (5)\\
Log-kill & fester Anteil je Gabe (13)\\
Lyapunow-Exponent & Rate des Fehlerwachstums (3)\\
Massenwirkung & Begegnungen $\propto x\cdot y$ (4)\\
Michaelis-Konstante & Halbsättigung (11.1)\\
Nernst-Spannung & Gleichgewicht je Ionenart (12)\\
Nullkline & Menge mit $\dot x = 0$ (5.2)\\
Phasenporträt & $y$ gegen $x$ statt über $t$ (4)\\
Quasi-Stationarität & schneller Vorgang eingestellt (11)\\
$\Rnull$ & Basisreproduktionszahl (7.2)\\
Refraktärzeit & Unerregbarkeit nach Ausschlag (12)\\
Rheobase & Mindestreizstärke (12)\\
Riegelungsimpfung & Impfen im Umfeld eines Falls (8)\\
Sattelpunkt & $D<0$, stets instabil (6)\\
Selektionsfenster & Spiegel wählt aus, ohne zu töten (14)\\
Strukturelle In\-stabilität & Verhalten bei kleiner Änderung weg (4)\\
Superspreading & wenige verursachen die meisten Fälle (9)\\
Überschwingen & Endgröße über der Schwelle (7)\\
Überstreuung & Varianz über dem Mittelwert (9)\\
Verteilungs\-volumen & $D/c_0$ (10.1)\\
Vorrangeffekt & wer zuerst da ist, bleibt (5)\\
Wendepunkt & größter Zuwachs (2.3, 13.3)\\
Zeitkonstante & $R_mC_m$ (12.1)\\
\bottomrule
\end{longtable}}
\end{center}
\end{merksatz}
